\chapter{Implementation}
\label{chap:implementation}
%%%%%%%%%%%%%%%%%%%%

This chapter describes the software and hardware implementation of the system designed in Chapter~\ref{chap:design}.

\section{Frontend and Telemetry Stack}
The user-facing layer is built within GSense, providing the orchestration dashboard, recording states, and test minigames. Telemetry ingestion is managed primarily by GXT-Devio via HIDPP and seamlessly synchronized by the Nightfall background service. This layered composition ensures that 1\,kHz hardware polling runs safely out-of-process, leaving the GSense frontend free to organize workflows without stalling the main analytical execution arrays.

\section{Machine Learning Backend}
The machine learning pipeline is encapsulated in a dedicated Python intent-detection backend. During preprocessing, SciPy-based MAKIMA interpolation and statistical variance filtering are executed to sanitize and score the 1\,kHz stream entropy. The Random Forest classifier relies heavily on scikit-learn logic, converting the sliding windows of distance/velocity probes into flattened tabular feature arrays. Experimental sequence-level architectures, specifically the FDCN approaches, are developed using PyTorch.

\section{Optimization Pipeline Integration}
The optimization phase bridges the machine learning framework back into direct hardware behavior suggestions. The CMA-ES simulation implements custom finite state machine logic algorithmically mirroring the precise internal trigger mechanisms for the PRO X2 Superstrike. The frontend integration directly extracts the optimizer's outputs to satisfy critical user-sentiment constraints observed during testing: displaying the resultant AP and RT adjustments not as a simple black-box recommendation, but as graph-supported "educated guesses" illustrating the explicit Recall/Precision tradeoff improvements.

\section{Data Acquisition Campaign}
\label{sec:data_acquisition}

The development and training of the machine learning backend required the construction of a robust ground-truth dataset spanning diverse ergonomic and biomechanical profiles. Data was collected from 69 participants via active field deployments at two major public gaming events (BIC and PolyLAN).

\subsection{Controlled Task Dataset (BIC Event)}
The first deployment cohort ($N=26$) comprised a mix of highly ranked competitive gamers and casual novices playing \textit{Counter-Strike 2} (CS2). The protocol ran for approximately 60 minutes per user and was strictly structured to isolate specific click dynamics:
\begin{itemize}
    \item \textbf{Configuration A/B Testing:} Following a uniform warmup phase, users played two 15-minute game blocks. The device hardware profile was hot-swapped halfway through the session, alternating between an aggressive configuration (AP at 1, RT at 1) and a highly conservative configuration (AP at 5, zero Rapid Trigger). The presentation order was randomized to mitigate sequence bias.
    \item \textbf{Weapon Rotations:} Within each 15-minute block, users were instructed to rotate their active weapon every five minutes. The arsenal was specifically chosen to evoke distinct motor patterns: the \textit{USPS} pistol for rapid individual tapping fatigue, the \textit{M4A1} rifle for sustained tracking and spraying, and the \textit{AWP} sniper rifle for slow, deliberate holding and release. 
\end{itemize}
Qualitative user feedback regarding the switch feel was collected asynchronously after the sessions via standardized surveying.

\subsection{Free-Play Dataset (PolyLAN Event)}
The secondary dataset ($N=43$) focused on unstructured free-play in a high-density LAN environment. Participants engaged in approximately 40-minute play sessions utilizing their preferred competitive title (\textit{League of Legends}, \textit{CS2}, or \textit{Valorant}). Similar to the controlled dataset, the hardware configurations were silently swapped from conservative to aggressive baseline profiles near the 20-minute mark to capture a wider net of organic responses.

The aggregated dataset generated from these two campaigns ultimately produced hundreds of hours of raw 1\,kHz telemetry, demanding the variance entropy scoring (detailed in Section~\ref{chap:design}) to isolate periods of genuine engagement before entering the classification pipelines.

\subsection{Annotated Dataset Summary}
Following manual annotation and variance-based trimming, the final curated dataset comprises 132 labeled recording files from 39 participants, totaling approximately 12,196 seconds (\textasciitilde3.4 hours) of verified active telemetry. Left-channel annotations account for 39,088 labeled click cycles and right-channel annotations for 6,645 cycles, reflecting the predominance of left-button interactions in competitive gaming. The dataset was partitioned into training, validation, and test splits performed at the participant level to enforce strict subject isolation.
